

\chapter{Grundlagen}

%%%%%

\section{Architektur und Syntaktik von MP4-Containern}

\begin{itemize}
    \item  Containerformat für MPEG-4-Inhalte
    \item  Seit 2001 standardisiert
    \item Basiert auf dem Apple-QuickTime-Dateiformat, das insbesondere für MOV-Dateien genutzt wurde
\end{itemize}

\subsection{Formale Spezifikation und hierarchische Baumstruktur}

\begin{itemize}
    \item Der MP4-Container basiert auf dem ISO Base Media File Format (ISO/IEC 14496-12) und ist in der Norm ISO/IEC 14496-14 standardisiert \cite{iso_14496-14_2020}
    \item Die Architektur ist objektorientiert und hierarchisch als Baumstruktur aufgebaut
    \item Die elementaren Datenstrukturen werden als Boxen oder Atome bezeichnet (z. B. ftyp, moov, mdat)
    \item Jede Box besteht aus einem Header, der die Länge und einen eindeutigen 4-Byte-Identifikator (Type) also durch einen Vier-Zeichen-Code (4CC) definiert, sowie den zugehörigen Daten oder weiter verschachtelten Sub-Boxen
    \item Herstellerspezifische Freiheitsgrade in der Anordnung und Nutzung optionaler Boxen (z. B. udta oder uuid) 
    \item Gerade diese optionalen Boxen werden von Herstellern und KI-Generatoren sehr unterschiedlich genutzt, und deren Existenz sowie Position sind besonders interessant
    \item TODO: Gerade ziehen der Bezeichnungen Box, Container, Atom, ...
\end{itemize}

\subsection{Funktionale Trennung von Metadaten und Bitstrom}

Struktur fordert Trennung zwischen Metadaten und dem eigentlichen audiovisuellen Bitstrom
ftyp-Box deklariert an erster Stelle die Kompatibilität und das Format der Datei
moov-Box fungiert als Top-Level-Container für alle Metadaten, Verknüpfungen und Timing-Informationen
trak-Boxen enthalten die Konfigurationen für jeden individuellen Medienstrom
mdat-Box speichert die rohen Audio- und Videodaten (für Container-Forensik nicht relevant)

\subsection{Herstellerspezifische Syntaktik und proprietäre Erweiterungen}



\begin{comment}
Definition des MP4-Containers als objektorientierte Baumstruktur. Die Datei besteht aus verschachtelten Boxen/Atomen

\begin{itemize}
    \item Hierarchische Baumstruktur gemäß ISO/IEC 14496-14
    \begin{itemize}
        \item Oft tief verschachtelte Objektstruktur
        \item Elementare Bausteine sind Atome/Boxen
        \item Box Typen (ftyp, moov, trak)
        \begin{itemize}
        \item ftyp: Definiert den Dateityp und grundlegende Kompatibilitätsmerkmale
        \item moov: Fungiert als Hauptcontainer für sämtliche Metadaten und die logische Struktur
        \item trak: Spureninformationen sowie udta, meta und hdlr für spezifische Metadaten
        \end{itemize}
        \item Metadaten (udta, meta, hdlr)
        \item Mediadaten (mdat -> uninteressant)
        \item Free Space (free, skip -> Anomalien? Vielleicht Interessant?)
    \end{itemize}
    \item Testboxen:
    \begin{itemize}
        \item udta (User Data Box): Container für benutzerdefinierte oder anwendungsspezifische Zusatzinformationen (z. B. Copyright-Tags, GPS-Koordinaten, Kameraeinstellungen
        \item colr (Colour Information Box): Spezifiziert den Farbraum, die Transfercharakteristika und die Farbmatrix des Videomaterials (z. B. ICC-Profile, BT.2020, nclx)
        \item free (Free Space Box): Ungenutzten Speicherplatz (Padding) innerhalb der Datei
        \item wide (Wide Box): Padding-Element (typischerweise 8 Byte groß), das historisch aus der Apple-QuickTime-Spezifikation stammt
        \item beam (Beam-/Loop-Box): WhatsApp verwendet diese, um Videos als GIFs darzustellen oder um sie für die Vorschau zu optimieren
    \end{itemize}
\end{itemize}

\subsection{Layout}

Jede Box folgt diesem binären Layout:

\begin{itemize}
    \item \textbf{Size (4 Bytes):} Gesamtlänge der Box (Big-Endian).
    \item \textbf{Type (4 Bytes):} Identifikator (z. B. \verb|ftyp|, \verb|moov|, \verb|mdat|).
    \item \textbf{Extended Size (Optional, 8 Bytes):} Falls Size = 1 (für Boxen > 4 GB).
    \item \textbf{Data:} Der eigentliche Inhalt oder weitere Sub-Boxen.
\end{itemize}

 
\begin{table}[H]
    \centering
\caption{Box-Typen aus der Testdatei \url{D01_V_flat_move_0001.mp4}}
\label{tab:placeholder}
    \begin{tabular}{|l|>{\raggedright\arraybackslash}p{0.85\linewidth}|}\hline
         \textbf{type}& \textbf{Beschreibung}
\\\hline
         ftyp& File Type Box: Deklariert Kompatibilität und Format (z. B. isom, 3gp4). Verifiziert die Datei als standardkonformes MP4.
\\\hline
         mdat& Media Data Box: Container für den gemultiplexten Binärstrom (Audio-/Video-Samples). Keine internen Metadaten.
\\\hline
         moov& Movie Box: Top-Level-Container für sämtliche Metadaten. Darf auf Dateiebene nur einmal vorkommen.
\\\hline
         mvhd& Movie Header Box: Enthält globale Timing-Parameter (Erstellungszeit, globale duration, timescale).
\\\hline
         trak& Track Box: Container für einen einzelnen, unabhängigen Medienstrom (z. B. Video oder Audio).
\\\hline
         tkhd& Track Header Box: Spezifiziert räumliche (Breite/Höhe in Festkommazahlen) und zeitliche Dimensionen des Tracks.
\\\hline
         mdia& Media Box: Container für die logische und zeitliche Definition der Track-Mediadaten.
\\\hline
         mdhd& Media Header Box: Definiert die trackspezifische Zeitbasis (timescale) und Dauer.
\\\hline
         hdlr& Handler Reference Box: Identifiziert den Medientyp des Tracks (z. B. vide für Video, soun für Audio).
\\ \hline
 minf&Media Information Box: Container für medienspezifische Metadaten (enthält u. a. stbl).
\\\hline
 stbl&Sample Table Box: Container für sämtliche Tabellen zur zeitlichen und physischen Lokalisierung von Samples.
\\\hline
 stsd&Sample Description Box: Definiert Codec-Spezifikationen (z. B. avc1, mp4a) und Initialisierungsparameter (SPS/PPS).
\\\hline
 stts&Time-to-Sample Box: Definiert die Dauer jedes Samples (Frames) in Ticks. Beweist VFR (Variable Framerate) oder CFR.
\\\hline
 stsz&Sample Size Box: Listet die exakte Größe (in Bytes) jedes einzelnen Samples im Track auf.
\\\hline
 stco&Chunk Offset Box: Listet die absoluten Byte-Offsets auf, bei denen Chunks (Gruppen von Samples) innerhalb der mdat beginnen.\\\hline
    \end{tabular}
    
    
\end{table}
\end{comment}


%%%%%



\section{Taxonomie des Hashing und Approximate Matching}
nutzen (insb. wiedererkennung)
Abgrenzung kryptografischer Hashfunktionen (kollisionsresistent, z. B. SHA-1) von Similarity Hashes (Approximate Matching)

\subsection{Kryptografische Hashfunktionen}

\subsection{Approximate Matching}

\subsection{Abstraktionsebenen}
abstraktionsebenen (byte, semantic, syntactic - das ist relevant

\subsection{Feature-Design}


%%%%%

\section{n-Gramm-Abstraktion und Set-Theorie}

 
%%%%%



\section{Distanzmetriken und Normalisierungsstrategien}

\subsection{Jaccard-Index (symmetrisch)}

\subsection{Tversky-Index (asymmetrisch, gewichtet)}

\subsection{Normalization Strategy}